\documentclass[11pt]{article}
\usepackage[margin=1in]{geometry}
\usepackage{amsmath,amssymb,amsthm,mathtools}
\usepackage{booktabs}
\usepackage{microtype}
\usepackage[hidelinks]{hyperref}
\usepackage{enumitem}

\newtheorem{theorem}{Theorem}
\newtheorem{lemma}[theorem]{Lemma}
\newtheorem{corollary}[theorem]{Corollary}
\newtheorem{proposition}[theorem]{Proposition}
\theoremstyle{definition}
\newtheorem{definition}[theorem]{Definition}
\newtheorem{remark}[theorem]{Remark}

\newcommand{\Z}{\mathbb{Z}}
\newcommand{\gr}{\operatorname{gr}}
\newcommand{\pd}{\operatorname{pd}}
\newcommand{\Hilb}{\operatorname{Hilb}}
\newcommand{\inI}{\operatorname{in}}
\newcommand{\vtwo}{\nu_2}

\title{\textbf{Collision Algebras and Exact Dyadic Ramification on the MPRC Ring}\\
Determinantal Presentation, Gröbner Normal Forms, and the $576$-Level Dense Local Algebra}
\author{Muhammad Arshad}
\date{September 2026}

\begin{document}
\maketitle

\begin{abstract}
We study the commutative algebra generated by the primitive movement collisions underlying the Top-KPQ geometry. Each sign orientation gives a quadratic Hankel determinantal ideal, and the two orientations share the sign-free zero and half-turn endpoints. We prove that the primitive $H(H-1)$ quadrics form a Gröbner basis with initial ideal
\[
(x_1,\ldots,x_{H-1})^2+(y_1,\ldots,y_{H-1})^2.
\]
The quotient is finite free of rank $H^2$ over the endpoint ring, with Hilbert series
\[
\frac{(1+(H-1)t)^2}{(1-t)^2},
\]
Krull dimension $2$, Cohen--Macaulay and Koszul structure, projective dimension $2H-2$, and a Betti polynomial equal to the square of the one-sign Eagon--Northcott polynomial.

We then specialize to the MPRC ring $\Z_{256}$ with $H=128$. On the dense diagonal locus, all relative differences collapse to one rotor $u$ with $u^{128}=1$. Writing $u=1+\varepsilon$, the nilpotent factor is
\[
T_8=\Z_{256}[\varepsilon]/((1+\varepsilon)^{128}-1).
\]
Using exact binomial valuations and the weight $W(2^a\varepsilon^b)=64a+b$, we prove
\[
\mathfrak m^{575}\ne0,\qquad \mathfrak m^{576}=0,
\]
with final nonzero layer generated by $128\varepsilon^{127}$. Hence the dense MPRC collision algebra has exact ramification length $576=9\cdot64$. A complete associated-graded presentation is left open.
\end{abstract}

\noindent\textbf{Keywords:} determinantal ideal; Hankel minors; Gröbner basis; Cohen--Macaulay; Koszul algebra; Eagon--Northcott; dyadic ramification; MPRC; $\Z_{256}$.

\section{Introduction}
The Top-KPQ completeness theorem gives a combinatorial normal form for the primitive even-cycle collision relations. This paper identifies the algebra generated by the same relations and then studies its dyadic specialization.

The algebraic route exposes two layers. Over a field, the collision relations form two coupled rational-normal-curve-type determinantal systems. Over the actual coefficient ring $\Z_{256}$, the dense diagonal locus is nonreduced: the two sign families coalesce through one relative rotor with a finite dyadic ramification tower.

The exact terminal depth is
\[
576.
\]
This number is derived internally from the $\Z_{256}$ relation and is not inserted from a vision patch size.

\paragraph{Contributions.}
\begin{enumerate}[leftmargin=*]
\item We identify each sign family with the $2\times2$ minors of a $2\times H$ Hankel matrix.
\item Using the collision normal form, we prove that the primitive quadrics are a Gröbner basis and compute the initial ideal.
\item We derive the free endpoint-module structure, Hilbert series, exact degree counts, Cohen--Macaulay and Koszul properties, and projective dimension.
\item We obtain the full Betti polynomial by tensoring the two one-sign Eagon--Northcott resolutions over the shared endpoint ring.
\item On the dense MPRC locus, we prove that all relative differences generate one parameter $u-1$ with $u^{128}=1$.
\item We prove the exact ramification index $576$ and identify the final nonzero layer.
\end{enumerate}

\paragraph{Novelty boundary.}
The field-side determinantal facts are classical: the one-sign ideal is the rational-normal-curve Hankel minor ideal, its Eagon--Northcott resolution is standard \cite{EagonNorthcott1962}, and combinatorial/chip-firing treatments of rational normal curves and their Gröbner degenerations already exist \cite{KarkiManjunath2023}. Likewise, once the dense factor is identified with the cyclic group ring $(\Z/2^8\Z)[C_{2^7}]$, the nilpotency index $576$ of the augmentation generator $u-1$ is a specialization of the classical cyclic augmentation-ideal formula; modern statements include \cite{Schauz2021}, building on earlier polynomial-periodicity results such as \cite{Wilson2006}. Our larger ideal is the maximal/Jacobson radical $(2,u-1)$. The classical augmentation formula alone does not establish that adjoining the coefficient radical $2$ leaves its nilpotency index unchanged; that is supplied here by the weight argument. We therefore claim no novelty for the general cyclic augmentation formula itself. The contribution asserted here is the MPRC-specific identification and synthesis, together with the exact full-radical calculation in this specialization.

\section{The collision algebra}
Let $k$ be a field and fix $H\ge2$. Introduce shared endpoint variables $z,h$, positive variables $x_1,\ldots,x_{H-1}$, and negative variables $y_1,\ldots,y_{H-1}$. Set
\[
x_0=y_0=z,\qquad x_H=y_H=h,
\]
and
\[
S=k[z,h,x_1,\ldots,x_{H-1},y_1,\ldots,y_{H-1}].
\]

Define
\[
I_+=
(x_i x_j-x_{i-1}x_{j+1}:1\le i\le j\le H-1),
\]
\[
I_-=
(y_i y_j-y_{i-1}y_{j+1}:1\le i\le j\le H-1),
\]
and
\[
I_H=I_++I_-.
\]
The collision algebra is
\[
A_H=S/I_H.
\]

\section{Determinantal form}
Consider
\[
M_x=
\begin{pmatrix}
z&x_1&x_2&\cdots&x_{H-1}\\
x_1&x_2&x_3&\cdots&h
\end{pmatrix}.
\]
The minor of columns $a<b$ equals
\[
x_a x_{b+1}-x_{a+1}x_b.
\]
Putting $i=a+1$ and $j=b$ gives exactly the primitive positive collision binomial. Hence $I_+$ is the ideal of $2\times2$ minors of $M_x$. The negative family is identical with $x$ replaced by $y$.

There are
\[
\binom H2=\frac{H(H-1)}2
\]
quadrics per sign and therefore
\[
\boxed{H(H-1)}
\]
primitive quadrics in total. At $H=128$ the count is $16{,}256$.

\section{Gröbner basis from collision confluence}
Define a variable weight
\[
w(z)=H^2,\qquad w(h)=0,\qquad
w(x_i)=w(y_i)=H^2-i^2\quad(1\le i\le H-1).
\]
Use a matrix monomial order whose first row is total degree, whose second row is the additive $w$-weight, and whose remaining rows give any lexicographic tie-breaker. The positive total-degree first row makes this a well-order. For every primitive source,
\[
w(x_i x_j)-w(x_{i-1}x_{j+1})
=2(j-i+1)>0,
\]
and identically for the $y$ variables. Thus
\[
x_i x_j
\]
is the leading monomial of
\[
x_i x_j-x_{i-1}x_{j+1}.
\]

Orient
\[
x_i x_j\longrightarrow x_{i-1}x_{j+1}
\]
and similarly for the $y$ family. The Top-KPQ collision theorem provides a strictly increasing squared-magnitude potential and a unique terminal for every collision class. Therefore the induced monomial rewriting is terminating and confluent for the displayed monomial order.

An irreducible monomial contains at most one ordinary positive interior variable and at most one ordinary negative interior variable. Thus the standard monomials are exactly
\[
z^a h^b q,
\qquad
q\in\{1,x_i,y_j,x_i y_j\}.
\]

\begin{theorem}[Quadratic Gröbner presentation]
For a term order compatible with the collision orientation,
\[
\boxed{
\inI(I_H)=
(x_1,\ldots,x_{H-1})^2+
(y_1,\ldots,y_{H-1})^2.
}
\]
The primitive $H(H-1)$ quadrics form a Gröbner basis.
\end{theorem}

\section{Free endpoint-module structure}
Let
\[
B=k[z,h].
\]
The standard monomials show that $A_H$ is a free $B$-module with basis
\[
\{1\}\cup\{x_i\}\cup\{y_j\}\cup\{x_i y_j\}.
\]
Its rank is
\[
1+2(H-1)+(H-1)^2=H^2.
\]

\begin{theorem}
\[
\boxed{
A_H\text{ is finite free of rank }H^2\text{ over }B=k[z,h].
}
\]
\end{theorem}

\section{Hilbert series and exact class count}
The $B$-basis has degree numerator
\[
1+2(H-1)t+(H-1)^2t^2=(1+(H-1)t)^2.
\]
Hence
\[
\boxed{
\Hilb(A_H,t)=\frac{(1+(H-1)t)^2}{(1-t)^2}.
}
\]

For movement degree $m\ge1$ (so a temporally ordered selection of $K$ positions has $m=K-1$), the coefficient is
\[
(m+1)+2(H-1)m+(H-1)^2(m-1),
\]
which simplifies to
\[
\boxed{
N_m=H[H(m-1)+2].
}
\]
This agrees with the independently derived Top-KPQ class count.

\section{Cohen--Macaulay, projective dimension, and Koszulness}
Because $A_H$ is finite free over $B=k[z,h]$, the endpoint variables form a regular system of parameters and
\[
\dim A_H=2.
\]
Thus $A_H$ is Cohen--Macaulay. Since $S$ has $2H$ variables,
\[
\boxed{
\pd_S A_H=2H-2.
}
\]
The defining Gröbner basis is quadratic, so the standard graded algebra is Koszul.

\section{Betti polynomial}
For one sign, let
\[
S_x=B[x_1,\ldots,x_{H-1}],\qquad R_x=S_x/I_+.
\]
The standard-monomial description shows that $R_x$ is finite free of rank $H$ over $B$, hence
\[
\dim R_x=2.
\]
Since $\dim S_x=H+1$, the minor ideal has height
\[
\operatorname{ht} I_+=H-1,
\]
which is the expected height of the maximal-minor ideal of a $2\times H$ matrix. Therefore the Eagon--Northcott complex is exact and gives the minimal resolution \cite{EagonNorthcott1962}. The quotient Betti numbers are
\[
\beta_i=i\binom H{i+1},
\qquad 1\le i\le H-1,
\]
with internal shift $i+1$.

Define
\[
B_H(u,t)
=
1+\sum_{i=1}^{H-1}
i\binom H{i+1}u^i t^{i+1}.
\]
Let
\[
R_x=B[x_1,\ldots,x_{H-1}]/I_+,
\qquad
R_y=B[y_1,\ldots,y_{H-1}]/I_-.
\]
Then
\[
A_H\cong R_x\otimes_B R_y.
\]
Both $R_x$ and $R_y$ are free, hence flat, over $B$. Let $F_x$ and $F_y$ be their Eagon--Northcott resolutions over the corresponding polynomial rings. The total complex of $F_x\otimes_B F_y$ is therefore exact and resolves $A_H$ over $S$: flatness eliminates higher $\operatorname{Tor}^B$, while every term is free over $S$. All differential entries have positive standard degree, so the total resolution remains minimal. Consequently
\[
\boxed{
B_{A_H}(u,t)=B_H(u,t)^2.
}
\]
The top homological degree is $2H-2$, matching the Cohen--Macaulay calculation.

\section{Dyadic specialization}
Set
\[
D=\Z/256\Z,\qquad H=128,
\]
and let $A_D$ be the collision algebra after base change to $D$. On the dense locus where $z$ and $h$ are units, the positive relations imply
\[
x_j=zt^j
\]
for a unit $t$, while the negative relations imply
\[
y_j=zs^j
\]
for a unit $s$. The shared endpoint gives
\[
zt^{128}=zs^{128},
\]
hence
\[
t^{128}=s^{128}.
\]

Define
\[
u=s/t.
\]
Then
\[
\boxed{u^{128}=1}
\]
and
\[
y_i=x_i u^i.
\]
Therefore
\[
y_i-x_i=x_i(u^i-1).
\]
Since $x_i$ is a unit and
\[
u^i-1=(u-1)(1+u+\cdots+u^{i-1}),
\]
all differences lie in $(u-1)$. Conversely,
\[
y_1-x_1=x_1(u-1),
\]
so
\[
\boxed{
(y_i-x_i:1\le i<128)=(u-1).
}
\]

Thus the dense diagonal ideal is
\[
\mathfrak p=(2,u-1).
\]

More precisely, if $A_{D,U}$ denotes the localization of the $D$-collision algebra at $zh$, then the above formulas give an explicit isomorphism
\[
\boxed{
A_{D,U}\cong
D[z^{\pm1},t^{\pm1}]
\otimes_D
D[u]/(u^{128}-1),
}
\]
via
\[
x_i\mapsto zt^i,\qquad
y_i\mapsto zt^iu^i,\qquad
h\mapsto zt^{128}.
\]
The inverse recovers $t=x_1/z$ and $u=y_1/x_1$. The Laurent polynomial factor
\[
R_U=D[z^{\pm1},t^{\pm1}]
\]
is faithfully flat over $D$. Hence powers and nonvanishing of the relative ideal are detected on the cyclic factor.

\section{The monic dense factor}
Put
\[
u=1+\varepsilon.
\]
Then
\[
(1+\varepsilon)^{128}=1.
\]
The nilpotent factor is
\[
T_8=
D[\varepsilon]/((1+\varepsilon)^{128}-1),
\qquad
\mathfrak m=(2,\varepsilon).
\]

The defining polynomial is monic of degree $128$, therefore
\[
1,\varepsilon,\ldots,\varepsilon^{127}
\]
is a free $D$-basis. In particular, divisibility by powers of $2$ is coefficientwise in this normal form.

\section{Exact binomial valuations}
For $1\le j<128$,
\[
j\binom{128}{j}=128\binom{127}{j-1}.
\]
Since $127=(1111111)_2$, every coefficient $\binom{127}{j-1}$ is odd. Hence
\[
\boxed{
\vtwo\binom{128}{j}=7-\vtwo(j).
}
\]
Thus
\[
\varepsilon^{128}
+
\sum_{j=1}^{127}
2^{7-\vtwo(j)}u_j\varepsilon^j=0,
\]
with every $u_j$ odd.

\section{The MPRC weight}
Define
\[
\boxed{
W(2^a\varepsilon^b)=64a+b.
}
\]
The leading monomial $\varepsilon^{128}$ has weight $128$. The $j$th lower term has weight
\[
64(7-\vtwo(j))+j.
\]
The unique lower term of the same weight occurs at $j=64$:
\[
W(2\varepsilon^{64})=64+64=128.
\]
Every other lower term has strictly larger weight.

Write
\[
\binom{128}{64}=2\omega,
\]
where $\omega\in D^\times$ is odd. Therefore the unique lowest-weight branch is
\[
\varepsilon^{128}
=
-2\omega\varepsilon^{64}
+
(\text{higher weight}).
\]
Every reduction is weight-nondecreasing.

\section{Exact ramification index}
A nonzero normal-form monomial has $0\le a\le7$ and $0\le b\le127$, hence
\[
W(2^a\varepsilon^b)\le7\cdot64+127=575.
\]

The ideal $\mathfrak m^{576}$ is generated by
\[
2^a\varepsilon^{576-a},\qquad 0\le a\le7.
\]
Their weights are
\[
64a+576-a=576+63a\ge576.
\]
Therefore every generator vanishes.

\begin{theorem}[Dense MPRC Ramification]
\[
\boxed{
\mathfrak m^{576}=0.
}
\]
\end{theorem}

To show sharpness, follow the unique equal-weight branch seven times:
\[
575\to511\to447\to383\to319\to255\to191\to127.
\]
This gives
\[
\varepsilon^{575}
=
(-2\omega)^7\varepsilon^{127}
+
(\text{weight}>575).
\]
The higher-weight terms vanish. Since $\omega$ is odd,
\[
(-2\omega)^7=-128\omega^7\equiv128\pmod{256}.
\]
Hence
\[
\boxed{
\varepsilon^{575}=128\varepsilon^{127}\ne0.
}
\]

For $a\ge1$, the remaining degree-$575$ generators satisfy
\[
W(2^a\varepsilon^{575-a})=575+63a>575,
\]
so they vanish. Therefore
\[
\mathfrak m^{575}=\langle128\varepsilon^{127}\rangle.
\]

\begin{corollary}[Exact index]
\[
\boxed{
\mathfrak m^{575}\ne0,\qquad
\mathfrak m^{576}=0,
\qquad
\mathfrak m^{575}=\langle128\varepsilon^{127}\rangle.
}
\]
Because $R_U$ is faithfully flat over $D$, the extended dense-diagonal ideal
\[
\mathfrak p=R_U\otimes_D\mathfrak m
\]
has the same exact nilpotency index in $A_{D,U}$.
\end{corollary}

\section{Relation to the classical cyclic augmentation ideal}
The factor
\[
T_8=D[u]/(u^{128}-1)
\]
is the group ring
\[
D[C_{128}].
\]
Its augmentation ideal is
\[
I_{\mathrm{aug}}=(u-1)=(\varepsilon),
\]
whereas
\[
\mathfrak m=(2,\varepsilon)
\]
is the larger maximal/Jacobson radical of $T_8$. Thus the classical augmentation result below applies directly to $(\varepsilon)$, not automatically to the full $\mathfrak m$. For
\[
(\Z/p^\beta\Z)[C_{p^\alpha}],
\]
the classical nilpotency index is
\[
\nu=(\beta(p-1)+1)p^{\alpha-1}
\]
in the cyclic case \cite{Schauz2021}; related coefficient bounds already appear in Wilson's periodic-polynomial theorem \cite{Wilson2006}. With
\[
p=2,\quad \alpha=7,\quad \beta=8,
\]
this gives
\[
\nu(I_{\mathrm{aug}})= (8+1)2^6=576.
\]
This independently confirms
\[
\varepsilon^{575}\ne0,\qquad \varepsilon^{576}=0.
\]
The preceding MPRC weight argument does more in the present setting: it proves that every mixed generator $2^a\varepsilon^{576-a}$ also vanishes, so the larger radical $\mathfrak m=(2,\varepsilon)$ has the same exact index $576$.

\section{The number $576$}
The proof yields the exact terminal boundary
\[
\boxed{
576=9\cdot64.
}
\]
Since $64=4\cdot16$,
\[
576=4\cdot9\cdot16=24^2.
\]
Only the identity $576=9\cdot64$ is presently internal to the ramification derivation. A structural identification with a $4\times4$ CXR observer or a $24\times24$ vision patch is not proved here.

For comparison, the established QH4 active-state count is
\[
252=4\cdot7\cdot9.
\]
The exact arithmetic transition
\[
(4,7,9)\to(4,9,16)
\]
is recorded as an open structural observation, not a theorem.

\section{Verification}
Independent integer reduction in the free $128$-element $D$-basis verifies
\[
\varepsilon^{575}=128\varepsilon^{127},\qquad
\varepsilon^{576}=0,
\]
and the vanishing of every degree-$576$ ideal generator. These checks are regression evidence; the proof is the monic-basis, valuation, and weight argument above.

\section{Open associated-graded problem}
The nilpotence theorem determines the endpoint of the filtration but not yet the complete presentation of
\[
\gr_{\mathfrak m}(T_8)
=
\bigoplus_{n\ge0}\mathfrak m^n/\mathfrak m^{n+1}.
\]
Exact finite calculations reveal a staircase, but this paper deliberately does not promote an empirical standard-basis pattern to a theorem. A complete analytical associated-graded presentation remains open.

\section{Conclusion}
The primitive MPRC collision relations define a rigid algebraic object. Over a field, they give two shared-endpoint Hankel determinantal systems with a quadratic Gröbner basis, exact Hilbert series, Cohen--Macaulay and Koszul structure, and squared Eagon--Northcott Betti data. Over $\Z_{256}$, their dense diagonal deformation collapses to one relative rotor $u^{128}=1$, whose dyadic local algebra has exact nilpotence index
\[
\boxed{576}.
\]
Thus the $576$ boundary is an intrinsic property of the dense collision algebra, not an imported architecture constant.

\section*{Reproducibility record}
The theorem register and freeze record are archived at
\begin{center}
\url{https://github.com/muhammadarshad/pages-mprc/tree/master/research/top-kpq}
\end{center}
and the exact collision/symmetry gates are archived at
\begin{center}
\url{https://github.com/muhammadarshad/siliq-rotor/tree/master/model/motif}.
\end{center}

\begin{thebibliography}{9}
\bibitem{TopKPQ}
M. Arshad,
\emph{Top-KPQ Completeness and Q-Space Geometry on Even Cyclic Rings},
preprint draft, 2026.

\bibitem{EagonNorthcott1962}
J. A. Eagon and D. G. Northcott,
\emph{Ideals defined by matrices and a certain complex associated with them},
Proceedings of the Royal Society of London. Series A,
269 (1962), 188--204.
doi:10.1098/rspa.1962.0170.

\bibitem{ArshadOperator}
M. Arshad,
\emph{An Eight-State Operator and Transpose Geometry over Dyadic Rings: Trace-Zero Modules, Mixed-Radix Addressing, and Tensor-Lift Kernels},
Zenodo, 2026.
doi:10.5281/zenodo.22850830.

\bibitem{KarkiManjunath2023}
R. Karki and M. Manjunath,
\emph{Rational Normal Curves, Chip Firing and Free Resolutions},
arXiv:2301.09104, 2023.

\bibitem{Schauz2021}
U. Schauz,
\emph{The Largest Possible Finite Degree of Functions between Commutative Groups},
arXiv:2103.16467, 2021.

\bibitem{Wilson2006}
R. M. Wilson,
\emph{A lemma on polynomials modulo $p^m$ and applications to coding theory},
Discrete Mathematics 306 (2006), 3154--3165.
doi:10.1016/j.disc.2004.10.030.
\end{thebibliography}

\end{document}
